Work in $\Cl_{\Q}(2,0)$, which exists by Theorem \claimref{synthesis.division-gives-the-rationals} (the integer algebra of Premise (F3) with every non-null vector inverted). Its even part $\{a + bB : a, b \in \Q\}$ is commutative, since $B$ commutes with $1$ and with itself, and $(a + bB)^{\sim} = a - bB$ because reversion reverses $e_{1}e_{2}$ to $e_{2}e_{1} = -B$. Hence $(a + bB)(a + bB)^{\sim} = a^{2} - b^{2}B^{2} = a^{2} + b^{2}$.

\emph{Invertibility.} $(1 + tB)(1 - tB) = 1 - t^{2}B^{2} = 1 + t^{2}$, a nonzero rational, so $(1 + tB)^{-1} = (1 - tB)/(1 + t^{2})$ and $R_{t} = (1 - tB)^{2}/(1 + t^{2}) = ((1 - t^{2}) - 2tB)/(1 + t^{2})$.

\emph{(i)} With $a = (1 - t^{2})/(1 + t^{2})$ and $b = -2t/(1 + t^{2})$, $a^{2} + b^{2} = ((1 - t^{2})^{2} + 4t^{2})/(1 + t^{2})^{2} = 1$, so $R_{t} \in U_{\Q}$ and $R_{t}\widetilde{R_{t}} = 1$. For $R \in U_{\Q}$ and $v \in V_{\Q}$, $RvR^{\sim}$ is a product of two even elements and one odd one, hence odd, and in $\Cl(2,0)$ the odd elements are the vectors, so $RvR^{\sim} \in V_{\Q}$; its square is $RvR^{\sim}RvR^{\sim} = Rv^{2}R^{\sim} = Q(v)$. The values $R_{0} = 1$ and $R_{1} = -2B/2 = -B$ are direct; $(-B)e_{1}(-B)^{\sim} = -Be_{1}B = -e_{1}$ and likewise $e_{2} \mapsto -e_{2}$, the half turn.

\emph{(ii)} The even part is commutative, so for rational $s, t$
\[
R_{s}R_{t} = \frac{(1 - sB)(1 - tB)}{(1 + sB)(1 + tB)} = \frac{(1 - st) - (s + t)B}{(1 - st) + (s + t)B}.
\]
If $st \ne 1$, divide numerator and denominator by $1 - st$: this is $R_{u}$ with $u = (s + t)/(1 - st)$. If $st = 1$ then $s + t \ne 0$ (as $s, t$ are rational with $st = 1 > 0$, they have the same sign), and the quotient is $(-(s + t)B)/((s + t)B) = -1 = R_{\infty}$. With infinity: $R_{\infty}R_{t} = -R_{t}$, and for $t \ne 0$, $R_{-1/t} = (t + B)/(t - B)$ (multiply numerator and denominator of $(1 + B/t)/(1 - B/t)$ by $t$), while $-R_{t} = (tB - 1)/(1 + tB)$; multiplying numerator and denominator of $(t + B)/(t - B)$ by $B$ gives $(tB - 1)/(tB + 1)$, so $-R_{t} = R_{-1/t}$. Also $R_{\infty}R_{0} = -1 = R_{\infty}$ and $R_{\infty}R_{\infty} = 1 = R_{0}$.

\emph{(iii)} We show the map is a bijection onto $U_{\Q}$; by (ii) it then carries $\oplus$ to the product of the group $U_{\Q}$, so $\oplus$ is a group law with identity $0$ and the map is an isomorphism. Injective: if $a + bB = R_{t}$ with $t$ rational then $1 + a = 2/(1 + t^{2}) \ne 0$ and $-b/(1 + a) = (2t/(1 + t^{2}))/(2/(1 + t^{2})) = t$, so $t$ is recovered, and $R_{t} \ne -1 = R_{\infty}$ since $1 + a \ne 0$. Surjective: let $a^{2} + b^{2} = 1$ with $a \ne -1$ and put $t = -b/(1 + a) \in \Q$. Then $1 + t^{2} = ((1 + a)^{2} + b^{2})/(1 + a)^{2} = (2 + 2a)/(1 + a)^{2} = 2/(1 + a)$, so $(1 - t^{2})/(1 + t^{2}) = (2 - (1 + t^{2}))(1 + a)/2 = a$ and $-2t/(1 + t^{2}) = b$; and $a = -1$ forces $b = 0$, the element $-1 = R_{\infty}$. Finally $R_{\infty}vR_{\infty}^{\sim} = (-1)v(-1) = v$.

\emph{(iv)} Over $\R$ (Premise (F4)) the formulas of (i) to (iii) hold verbatim for real $t$, the even part with $|a + bB|^{2} = a^{2} + b^{2}$ being the Euclidean plane in the coordinates $(a, b)$. For $t \ge 0$, $b(t) = -2t/(1 + t^{2}) \le 0$, and $a(t) = (1 - t^{2})/(1 + t^{2})$ has $a'(t) = -4t/(1 + t^{2})^{2} < 0$ for $t > 0$, so $a$ decreases strictly from $1$ at $t = 0$ towards $-1$ as $t \to \infty$; since $a^{2} + b^{2} = 1$ and $b \le 0$, the point $(a, b)$ is determined by $a$, and the curve runs once along the lower half of the unit circle from $(1, 0)$ to $(-1, 0)$, the endpoint being $R_{\infty}$. Its length is half the length of the unit circle, which is $\pi$ by definition. Differentiating, $b'(t) = -2(1 - t^{2})/(1 + t^{2})^{2}$, and
\[
a'(t)^{2} + b'(t)^{2} = \frac{16t^{2} + 4(1 - t^{2})^{2}}{(1 + t^{2})^{4}} = \frac{4(1 + t^{2})^{2}}{(1 + t^{2})^{4}} = \frac{4}{(1 + t^{2})^{2}},
\]
so the speed is $2/(1 + t^{2})$ and the length is $\int_{0}^{\infty} 2\,dt/(1 + t^{2}) = \pi$. $\square$

The check verifies (i) to (iii) in exact rational arithmetic with GA.py (the rotor property and the vector action for a grid of rationals, the composition law on all pairs from a set of rationals including the cases $st = 1$ and infinity, and the bijection on every rational point of the circle with height at most $60$), the speed identity as an identity of rational functions, and brackets the length $\int_{0}^{\infty} 2\,dt/(1 + t^{2})$ between rational Riemann-sum bounds, which it compares with rational bounds on $\pi$ from Machin's formula.
